\subsection{Conclusions}
\label{sec:conclusions-answers}

Wu et al.\ priced their physics-informed effective Hamiltonian by the gates in one query and left open how its cost depends on the spectral gap and on the cost of differentiation \cite{wu2025pihm}. Restated as the standard form, whose state holds the field alone and represents every derivative through the dense Chebyshev differentiation matrix $\G$, their construction reproduces the operators they printed and the scale factors of their ordinary examples once the printed choices are restored. Those choices (rounded constraint points, truncated sources and invariant sets read off the known solution) leave error floors or a degenerate kernel that no increase in resolution cures, and the standard form replaces them; their partial differential and nonlinear examples were not reproduced. Question (i) of \Cref{sec:pihm-assessment} asked how the number of queries, the filter's degree, grows with the resolution, and how much of that growth is the encoder's and how much the Hamiltonian's. With Wu et al.'s approximate encoding of $\G$, most of the degree is the encoder's, a looseness that grows with $N$ and that an exact encoding of the same Hamiltonian removes. What remains is the Hamiltonian's: in the standard form, for an order-$k$ equation with constant coefficients, the degree is $\Omega(N^{2k}\log(1/\varepsilon))$ under any block encoding and every polynomial filter of the walk, and it grows as $\tilde\Theta(N^{2k})$ with the structured encoding, where the gap is bounded above by proof and measured not to fall. The filter sets only the constant. The minimax edge filter, optimal for the worst case that the spectral edge alone determines, needs a lower degree than imaginary time at equal suppression, and neither filter changes how the degree grows with $N$.

Question (ii) asked whether the residual can be rewritten so that its Hamiltonian is banded and well conditioned without changing its solution, and at what cost. The descriptor form keeps apart the two banded factors whose product is $\G$, carrying a higher-order equation's derivatives as unknowns and multiplying a first-order equation or system through by the conversion factor. Its kernel is the standard form's. On a regular problem with one axis, polynomial coefficients, a leading coefficient without zeros on $[-1,1]$ and analytic data, its degree grows as $\tilde\Theta(N)$ independent of order. The conditioning moves the norm's cost to the gap. On a singular problem whose leading coefficient has a simple zero, block rescaling restores a gap of order one on a single axis, without ancilla and at a cost in gates fixed by the equation, for a degree of about $\tilde\Theta(N^2)$. On two axes the gap falls with $N$ whatever the weights. The descriptor's other prices are a register of derivative blocks, whose size the equation fixes, and, on constant-coefficient problems at small sizes, more gates per query than the structured standard form, a difference that narrows as the resolution grows, so its advantage lies wholly in the degree.

Question (iii) asked whether the solution of a nonlinear equation can be prepared as a unique ground state, and whether the reformulation's advantage survives there. Wu et al.'s doubled space encodes a polynomial nonlinearity exactly, but its zero-energy eigenspace has dimension at least $2^{2n-1}$, so a filter can only project onto it. A truncated Carleman lift has an isolated ground state, which both forms prepare, from Burgers' equation to two-dimensional incompressible Navier--Stokes in exact emulation and resource estimate. The descriptor's advantage survives through the time axis --- which the lift still differentiates --- once that axis is resolved finely enough. The uniqueness that makes preparation possible certifies the truncated problem only; a lift past its boundary keeps an isolated ground state while its field departs from the nonlinear solution. The condition $\rho<1$ under which Liu et al.\ prove that the truncation converges is a worst case over all data, which structured data need not meet \cite{liu2021carleman}.

The three answers rest on one pipeline shared by both forms, whose circuits were verified against the operators they encode and whose exact emulation equals the circuit's post-selected output wherever both ran, so a difference measured between the forms is a difference between their residuals. They describe one mechanism: the cost of preparing the ground state of a differential equation is set by the conditioning of its residual, which the formulation fixes before any circuit is encoded. Neither form, on a linear equation or a lifted one, is polylogarithmic in $N$. Because built circuits favour the form with the tighter encoding, a comparison speaks about the Hamiltonians only at the ideal-encoding bound. The degree has the form of the optimal law for linear systems, with the residual's effective condition number in place of the matrix's, so the ground-state route meets conditioning in another guise. The tools for lowering it come from classical numerical analysis, two of which the descriptor transferred: the banded operators of the ultraspherical method \cite{olver2013} and the carried derivatives of first-order-system least squares \cite{cai1994fosls}.

No speed-up over classical solvers is claimed. An analytic solution needs a resolution growing only as $\log(1/\delta)$ for an accuracy $\delta$, and on one axis a degree linear in $N$ matches the ultraspherical method's classical operation count. What the reformulation achieves is a preparation whose degree tracks the classical cost in order on a regular problem, and a route by which nonlinear flows are prepared as unique ground states while their lift converges. An advantage could appear only on several axes, where a classical solver's cost grows at least as the number of unknowns, and that is where the descriptor's gap still falls, so whether a residual can keep its gap on several axes decides whether this route can compete with classical spectral solvers at all.

\subsection{Future Work}
\label{sec:future-work}

On several axes the obvious repair, weighting the datum rows by a power of the tangential mode, has been measured: it raises the gap more than the norm and lowers the cost modestly, but the weighted gap still falls with every qubit. The question is now whether any reweighting, alone or with block rescaling, or rows of the kind first-order-system least squares adds to make its functional coercive, can make the residual bound every carried derivative on several axes. The most direct such rows impose the boundary data on the carried derivative blocks, through the data's tangential derivatives, as well as on the field. The solution satisfies them, so they would leave an exact kernel unchanged, and the gap and norm they give on Laplace's equation (R5a) at three to five qubits per axis would test the idea directly. The other candidate is the full ultraspherical reduction (\Cref{sec:limitations}), a banded residual on the field alone that carries no derivative. Building it and measuring one problem of each class, R2a, Legendre's equation (R3) and the heat equation (R5b), would show which classes it serves.

Block rescaling sets its weight from the equation's coefficients and the resolution, but its two constants were selected by calibration on the benchmark problems and then held fixed, so its check is in-sample. The quotient of the states concentrated near a zero already fixes how $\zeta$ must grow with $N$, so deriving the weight from the equation, its constants and the crossover between the classes included, should be possible. A regular coefficient that dips towards zero, where the classes meet, is where a derived weight would first be tested against the best on a grid. Block rescaling weights the kernel by one constant per block and so stays cheap, whereas an equilibration per amplitude would need a block encoding of its own at a cost linear in the dimension. Whether a diagonal structured in the Chebyshev index, cheap to encode where the general one is not, can do for the degree what diagonal preconditioning does for the ultraspherical method's banded systems \cite{olver2013}, or whether a degree linear in $N$ is the least that a differentiating residual reaches on one axis, remains open.

Every preparation here also takes the spectral edge from the classical reference. An underestimated gap keeps the target at a degree that grows as the inverse square root of the estimate (\Cref{app:edge-study}), so a lower bound loose by a constant costs only a constant, and what remains open is how to obtain such a bound without solving the problem classically. Lin and Tong estimate the ground energy on the same kind of block encoding, but their preparation still assumes a lower bound on the gap \cite{lintong2020groundstate}. On a regular problem \Cref{thm:flat-gap} proves the descriptor's gap bounded below but, by compactness, gives no constant. An explicit constant in terms of the equation's coefficients would supply the bound before any run, and would also fix how far the gap can fall on a stiff equation such as R2c.

Any condition that annihilates the solution enters the residual as rows, which is how coupled fields and realified complex-valued equations reached the same filter. A complex residual still gives a Hermitian Gram operator, so native complex preparation needs complex loaders and no new filter, and $z'=-\mathrm{i}\pi z$, which R7 carries realified, is the natural first comparison of the two encodings. An observation $f(x_i)=y_i$ would enter as the row $y_sf(x_i)-y_if(x_s)=0$. Adding observations of growing noise to a regular problem such as R2a, and following the ground value and the gap, would show how far a physics-informed Hamiltonian can become a data-informed one. Both directions change the rows and leave the filter and its verification as they are.

Nothing here ran on hardware, and the method's robustness to noise is unknown. A filter applies its queries in sequence, so an error in each query accumulates over the walk, and the error-corrected machine that the estimates assume must hold its logical error per gate well below the inverse of the preparation's gate count. Simulating the smallest preparation, the decay R1 at $n=2$, under a depolarising noise model would measure how the prepared state's infidelity grows with the degree and how much of the error the post-selection on the filter's ancilla discards. Combined with the estimated gate counts, it would set the logical error rate that each problem needs.
